\section{The system}

Figure \ref{fig:system} shows the whole system.
It has five parts, each drawn in its own colour in this and the following figures.

\begin{figure}[p]
\centering
% Figure 5.2: the whole system in place. Coordinates in cm; the page is 12.6 cm wide.
\begin{tikzpicture}[ggm, x=0.99cm, y=1cm]
  \newcommand{\frametitle}[3]{\node[font=\sffamily\scriptsize\bfseries, text=#1, anchor=north west, inner sep=1.5mm] at (#2) {#3};}
  \newcommand{\tinylabel}[2]{\node[font=\sffamily\tiny, text=ggmInk, align=center, anchor=north] at (#1) {#2};}

  % ================= external sources (top) =================
  \filldraw[draw=ggmSource!70, fill=ggmSource!7, rounded corners=3mm] (0,14.6) rectangle (12.5,16.6);
  \frametitle{ggmSource!80!black}{0,16.6}{EXTERNAL SOURCES}
  \foreach \x/\t in {1.55/{ENTSOG, GIE}, 3.75/{EIA}, 5.95/{Global Energy\\Monitor}, 8.15/{Outlooks and\\statistics}, 10.6/{Country studies\\and literature}} {
    \pic[scale=0.75] at (\x,15.75) {doc};
    \tinylabel{\x,15.3}{\t}
  }

  % ================= analyst's workspace (left) =================
  \filldraw[draw=ggmAnalyst!60, fill=ggmAnalyst!5, rounded corners=3mm] (0,7.9) rectangle (6.75,14.1);
  \frametitle{ggmAnalyst!85!black}{0,14.1}{THE ANALYST'S WORKSPACE}
  % floor shadow
  \fill[ggmModel!10] (2.9,8.15) ellipse (2.6 and 0.12);
  % chair
  \fill[ggmModel!40, rounded corners=1mm] (0.35,9.0) rectangle (0.55,10.6);
  \fill[ggmModel!40, rounded corners=1mm] (0.35,8.95) rectangle (1.45,9.12);
  \draw[ggmModel!55, line width=1.4pt] (0.9,8.95) -- (0.9,8.35);
  \draw[ggmModel!55, line width=1.4pt] (0.45,8.3) -- (1.35,8.3);
  % legs
  \draw[ggmModel!85, line width=5pt] (0.95,9.18) -- (1.75,9.18) -- (1.75,8.35);
  \fill[ggmInk, rounded corners=0.5mm] (1.65,8.22) rectangle (2.05,8.36);
  % torso
  \fill[ggmAnalyst] (0.6,9.12) .. controls (0.45,10.0) and (0.65,10.55) .. (1.0,10.65)
     .. controls (1.35,10.55) and (1.45,9.9) .. (1.35,9.12) -- cycle;
  % arm towards the keyboard
  \draw[ggmAnalyst!80!black, line width=4.2pt] (1.1,10.25) -- (1.45,9.75) -- (2.05,9.82);
  \fill[ggmSkin] (2.1,9.82) circle (0.075);
  % head
  \fill[ggmSkin] (1.05,10.98) circle (0.25);
  \fill[ggmHair] (0.79,11.03) .. controls (0.8,11.33) and (1.2,11.38) .. (1.3,11.12)
     .. controls (1.15,11.14) and (1.0,11.08) .. (0.97,10.97) .. controls (0.9,10.88) and (0.82,10.9) .. (0.79,11.03);
  % desk
  \fill[ggmSource!35!white] (1.55,9.62) rectangle (6.55,9.76);
  \fill[ggmSource!55!black] (1.75,8.3) rectangle (1.85,9.62);
  \fill[ggmSource!55!black] (6.25,8.3) rectangle (6.35,9.62);
  \fill[ggmModel!50, rounded corners=0.4mm] (1.95,9.76) rectangle (2.85,9.83);
  \fill[ggmModel!50] (3.05,9.76) ellipse (0.09 and 0.05);
  % coffee mug
  \fill[ggmAnalyst!70] (5.95,9.76) rectangle (6.17,10.0);
  \draw[ggmAnalyst!70, line width=0.8pt] (6.17,9.92) arc (90:-90:0.06);
  % monitor
  \fill[ggmModel!70] (4.0,9.76) -- (4.5,9.76) -- (4.38,10.0) -- (4.12,10.0) -- cycle;
  \fill[ggmModel, rounded corners=1.2mm] (2.0,9.98) rectangle (6.45,13.55);
  \fill[white] (2.1,10.1) rectangle (6.35,13.45);
  % window bar
  \fill[ggmModel!12] (2.1,13.2) rectangle (6.35,13.45);
  \foreach \i/\c in {0/ggmAnalyst, 1/ggmSource, 2/ggmRecord} \fill[\c] (2.22+0.14*\i,13.325) circle (0.045);
  \node[font=\sffamily\tiny, text=ggmModel] at (4.22,13.325) {GGM workbench};
  % spreadsheet pane
  \fill[ggmSheet] (2.18,12.88) rectangle (4.18,13.12);
  \node[font=\sffamily\tiny\bfseries, text=white, anchor=west] at (2.2,13.0) {Sheet: Arcs};
  \foreach \r in {0,...,9} {
    \pgfmathsetmacro{\y}{12.88-0.26*(\r+1)}
    \ifnum\r=4 \fill[ggmAgent!18] (2.18,\y) rectangle (4.18,\y+0.26); \fi
    \draw[ggmSheet!25] (2.18,\y) -- (4.18,\y);
    \foreach \c/\w in {0/0.32, 1/0.42, 2/0.28, 3/0.3, 4/0.26} {
      \pgfmathsetmacro{\x}{2.24+0.4*\c}
      \pgfmathsetmacro{\ww}{\w*(0.75+0.25*mod(\r*7+\c*3,4)/3)}
      \fill[ggmModel!25, rounded corners=0.3mm] (\x,\y+0.09) rectangle (\x+\ww,\y+0.17);
    }
  }
  \foreach \c in {1,...,4} \draw[ggmSheet!25] (2.18+0.4*\c,10.28) -- (2.18+0.4*\c,12.88);
  \draw[ggmSheet!40] (2.18,10.28) rectangle (4.18,12.88);
  \fill[ggmAgent] (4.12,11.71) circle (0.06);
  % suggestion card
  \filldraw[draw=ggmAgent!50, fill=ggmAgent!4, rounded corners=0.8mm] (4.32,10.2) rectangle (6.27,13.1);
  \fill[ggmAgent, rounded corners=0.8mm] (4.32,12.8) rectangle (6.27,13.1);
  \fill[ggmAgent] (4.32,12.8) rectangle (6.27,12.9);
  \pic[scale=0.22] at (4.47,12.95) {agent badge};
  \node[font=\sffamily\tiny\bfseries, text=white, anchor=west] at (4.55,12.95) {Suggestion};
  \node[font=\sffamily\tiny\bfseries, text=ggmInk, anchor=west] at (4.36,12.6) {New arc capacity};
  \node[font=\sffamily\tiny, text=ggmMuted, anchor=west] at (4.36,12.33) {Why};
  \foreach \y/\w in {12.12/1.75, 11.95/1.5, 11.78/1.65} \fill[ggmModel!25, rounded corners=0.3mm] (4.4,\y) rectangle (4.4+\w,\y+0.08);
  \node[font=\sffamily\tiny, text=ggmMuted, anchor=west] at (4.36,11.5) {Source};
  \fill[ggmModel!25, rounded corners=0.3mm] (4.4,11.25) rectangle (5.9,11.33);
  \node[fill=ggmRecord, text=white, rounded corners=0.6mm, font=\sffamily\tiny, inner xsep=1mm, inner ysep=0.4mm, anchor=west] at (4.38,10.95) {official, recent};
  \node[fill=ggmAnalyst, text=white, rounded corners=0.6mm, font=\sffamily\tiny\bfseries, inner xsep=0.5mm, inner ysep=0.5mm] at (4.66,10.47) {Accept};
  \node[draw=ggmModel!50, text=ggmModel, rounded corners=0.6mm, font=\sffamily\tiny, inner xsep=0.5mm, inner ysep=0.5mm] at (5.3,10.47) {Edit};
  \node[draw=ggmModel!50, text=ggmModel, rounded corners=0.6mm, font=\sffamily\tiny, inner xsep=0.5mm, inner ysep=0.5mm] at (5.87,10.47) {Reply};

  % ================= agents (right) =================
  \filldraw[draw=ggmAgent!60, fill=ggmAgent!5, rounded corners=3mm] (7.85,7.9) rectangle (12.5,14.1);
  \frametitle{ggmAgent}{7.85,14.1}{AGENTS}
  \node[font=\sffamily\tiny, text=ggmMuted, anchor=north east, inner sep=1.5mm] at (12.5,14.1) {one language model, four roles};
  \foreach \i/\n/\d in {0/{Scenario agent}/{storyline $\to$ edits},
                        1/{Data agent}/{sources $\to$ sourced values},
                        2/{Calibration agent}/{deviations $\to$ changes},
                        3/{Results agent}/{results $\to$ explanations}} {
    \pgfmathsetmacro{\y}{12.75-1.25*\i}
    \filldraw[draw=ggmAgent!40, fill=white, rounded corners=2mm, ggm soft shadow] (8.1,\y-0.48) rectangle (12.35,\y+0.48);
    \pic[scale=0.62] at (8.62,\y) {agent badge};
    \node[font=\sffamily\scriptsize\bfseries, text=ggmInk, anchor=west] at (9.05,\y+0.17) {\n};
    \node[font=\sffamily\tiny, text=ggmMuted, anchor=west] at (9.05,\y-0.17) {\d};
  }

  % suggestion and reply arrows between workspace and agents
  \draw[ggm flow=ggmAgent, line width=1.1pt] (7.95,12.6) -- (6.5,12.6);
  \node[font=\sffamily\tiny, text=ggmAgent, align=center, anchor=south] at (7.2,12.65) {suggestion\\+ reasons};
  \draw[ggm flow=ggmAnalyst, line width=1.1pt] (6.5,11.3) -- (7.95,11.3);
  \node[font=\sffamily\tiny, text=ggmAnalyst!85!black, align=center, anchor=north] at (7.2,11.25) {accept, edit\\or reply};

  % sources to agents
  \draw[ggm flow=ggmSource!80!black] (10.2,14.6) -- (10.2,14.15);
  \node[font=\sffamily\tiny, text=ggmSource!80!black, anchor=west] at (10.25,14.38) {fetch and read};

  % ================= tools (middle band) =================
  \filldraw[draw=ggmCode!60, fill=ggmCode!6, rounded corners=3mm] (0.75,5.05) rectangle (12.5,7.55);
  \frametitle{ggmCode}{0.75,7.55}{TOOLS: ORDINARY CODE}
  \foreach \i/\t in {0/{Read and write\\workbooks}, 1/{Convert and\\aggregate}, 2/{Shipping\\distances}, 3/{Consistency\\checks},
                     4/{Compare\\versions}, 5/{Deviation\\report}, 6/{Apply changes\\within bounds}, 7/{Start a\\model run}} {
    \pgfmathsetmacro{\x}{1.0+(mod(\i,4))*2.775+((mod(\i,4)>1) ? 0.45 : 0)}
    \pgfmathsetmacro{\y}{(\i<4) ? 6.25 : 5.25}
    \filldraw[draw=ggmCode!35, fill=white, rounded corners=1.5mm] (\x,\y) rectangle (\x+2.575,\y+0.78);
    \pic[scale=0.48] at (\x+0.35,\y+0.39) {gear};
    \node[font=\sffamily\tiny, text=ggmInk, align=left, anchor=west] at (\x+0.68,\y+0.39) {\t};
  }
  % agents call tools
  \draw[ggm flow=ggmAgent] (10.2,7.9) -- (10.2,7.6);
  \node[font=\sffamily\tiny, text=ggmAgent, anchor=west] at (10.25,7.75) {act only through tools};

  % ================= shared records (bottom left) =================
  \filldraw[draw=ggmRecord!60, fill=ggmRecord!6, rounded corners=3mm] (0.75,0.95) rectangle (7.6,4.65);
  \frametitle{ggmRecord}{0.75,4.65}{SHARED RECORDS}
  \foreach \x/\t in {1.55/gen, 2.45/proj, 3.35/calib} \pic[scale=0.95] at (\x,2.95) {workbook=\t};
  \tinylabel{2.45,2.35}{The three workbooks,\\under version control}
  \pic[scale=0.95] at (4.75,2.95) {database};
  \tinylabel{4.75,2.35}{Provenance\\ledger}
  \pic[scale=1.0] at (5.9,2.95) {doc=ggmRecord};
  \tinylabel{5.9,2.35}{Lessons}
  \pic[scale=0.95] at (6.95,2.95) {database=ggmModel};
  \tinylabel{6.95,2.35}{Run log}
  \node[font=\sffamily\tiny, text=ggmMuted, align=center] at (4.2,1.3) {every value: where it came from, who proposed it, who accepted it};
  % tools write to records
  \draw[ggm flow=ggmCode] (2.45,5.05) -- (2.45,3.55);
  \draw[ggm flow=ggmCode] (4.75,5.05) -- (4.75,3.55);

  % ================= local machine (bottom right) =================
  \filldraw[draw=ggmModel!60, fill=ggmModel!6, rounded corners=3mm] (8.1,0.95) rectangle (12.5,4.65);
  \frametitle{ggmModel}{8.1,4.65}{LOCAL MACHINE}
  \pic[scale=1.75] at (10.35,2.35) {laptop};
  \node[font=\sffamily\scriptsize\bfseries, text=ggmModel] at (10.35,3.1) {GGM};
  \node[font=\sffamily\tiny, text=ggmModel] at (10.35,2.77) {Julia + Gurobi};
  \node[font=\sffamily\tiny, text=ggmMuted, align=center] at (10.35,1.4) {model and licensed data\\stay on the analysts' machine};
  \draw[ggm flow=ggmModel] (9.55,5.05) -- (9.55,3.6);
  \node[font=\sffamily\tiny, text=ggmModel, anchor=east] at (9.5,3.85) {start run};
  \draw[ggm flow=ggmModel] (11.15,3.6) -- (11.15,5.05);
  \node[font=\sffamily\tiny, text=ggmModel, anchor=west] at (11.2,3.85) {results only};

  % ================= the analyst's decisions and the lessons =================
  \draw[ggm flow=ggmAnalyst, line width=1.1pt, rounded corners=3mm] (0.35,8.6) -- (0.35,3.9) -- (0.72,3.9);
  \node[font=\sffamily\tiny, text=ggmAnalyst!85!black, rotate=90, anchor=south] at (0.3,6.3) {decisions and corrections};
  \draw[ggm flow=ggmRecord, line width=1.1pt, dashed, rounded corners=3mm] (6.05,3.5) -- (6.675,4.3) -- (6.675,7.72) -- (7.3,7.72) -- (7.3,8.45) -- (7.82,8.45);
  \node[font=\sffamily\tiny\bfseries, text=ggmRecord, fill=ggmCode!6, inner sep=0.4mm, rotate=90] at (6.675,6.3) {lessons};
\end{tikzpicture}

\caption{The proposed system at work. The analyst works on the workbooks as today, with a panel of suggestions from four agents. The agents use tools written as ordinary code, read from external sources, and write to shared records. GGM runs on the analysts' own machine, and only its results return to the system.}
\label{fig:system}
\end{figure}

\textbf{The analyst's workspace.} The analyst works on the three workbooks, as today, but next to them sits a panel of suggestions.
Each suggestion is a proposed piece of work, such as a set of capacities for new nodes or a change to a production cost, shown with its reasons, sources and confidence.
The analyst accepts it, edits it, or replies in plain language.

\textbf{Four agents.} Each covers part of the workflow.
A scenario agent turns a storyline into proposed edits and a node structure.
A data agent fetches and extracts source data, maps it to nodes, and drafts the projections and demand parameters.
A calibration agent diagnoses deviations and proposes parameter changes.
A results agent compares runs and drafts explanations.
All four can share one underlying language model; they differ in their instructions, their tools and what they are allowed to change.

\textbf{Tools written as ordinary code.} The agents act only through tools: functions that read and write workbooks, convert units, aggregate capacities, compute shipping distances, check the workbooks for consistency, compare versions, compute deviation reports and start model runs.
These are where every number is computed, and they can be tested once and trusted afterwards.

\textbf{Shared records.} Four records hold the state of a study.
The workbooks themselves are kept under version control, so every change can be traced and undone.
A provenance ledger records, for every value, where it came from, who proposed it and who accepted it.
A set of lessons holds the general rules drawn from the analyst's corrections.
A run log records which workbooks and which version of the code produced each result.

\textbf{The model on a local machine.} GGM and its solver run where they run today.
The system starts runs and reads their results, but sends neither the model nor its licensed inputs elsewhere.

Most of this is ordinary software.
The language model appears only where the analysis of Section 5.2 found reading, mapping, diagnosis or drafting, and everything it produces passes through a tool, a record and, where it matters, the analyst.
